\chapter{Variations, extrémums et optimisation}
\section*{Exercice 16.1 -- Lire des variations}
Une fonction décroît de $8$ à $-1$ sur $[-2;1]$, puis croît de $-1$ à $4$ sur $[1;5]$. Donner ses extrémums sur $[-2;5]$. Peut-on comparer $f(0)$ et $f(4)$ ?

\section*{Exercice 16.2 -- Borne non atteinte}
La fonction $f(x)=x^2$ possède-t-elle un maximum et un minimum sur $]0;2[$ ?

\section*{Exercice 16.3 -- Optimiser}
Un rectangle a pour périmètre $24$ mètres. Déterminer son aire maximale en justifiant qu'elle est atteinte.

\section*{Exercice 16.4 -- Dichotomie}
Combien d'étapes de dichotomie depuis un intervalle d'amplitude $1$ garantissent une amplitude au plus $0{,}01$ ?

\section*{Exercice 16.5 -- Grille et preuve}
Un programme cherche le maximum de $f(x)=1-(x-0{,}3)^2$ sur la grille $\{0,0{,}5,1\}$. Quel résultat trouve-t-il ? Quel est le véritable maximum sur $[0;1]$ ?

\section*{Exercice 16.6 -- Longueur approchée}
Pour $y=x^2$ sur $[0;1]$, calculer la longueur polygonale avec un seul segment, puis avec les abscisses $0,1/2,1$.
